% Internal anonymous LNCS working draft, 8 October 2026.
% Pending result cells must be filled from verified originals.
\documentclass[runningheads]{llncs}
\usepackage[T1]{fontenc}
\usepackage{amsmath,amssymb,booktabs}
\newcommand{\method}{TitanTPP}
\newcommand{\softplus}{\operatorname{softplus}}
\newcommand{\Nrm}{\operatorname{N}}
\begin{document}
\title{Explicit Local and Recurrent History Correction for Event Quantity Prediction}
\titlerunning{Explicit History Correction for Event Quantities}
\author{Anonymous}
\authorrunning{Anonymous}
\institute{Anonymous institution}
\maketitle

\begin{abstract}
An event encoder can summarize earlier observations, but this does not
establish whether a separate path from recent observations improves
quantity prediction. We examine this question in a two-block causal
encoder for next-event quantity and recorded-time prediction.
A local three-position correction modifies the first block's query,
key, and value projections; a recurrent residual connects the first
and second blocks. Controls remove this path or replace its
cross-position inputs with current-position operations at the same
active parameter count. A separate study tests fixed combination
of time- and quantity-conditioned prototype scores against uniform,
similarity-weighted, and parameter-matched general-query retrieval.
The protocol fixes observable histories, output heads, checkpoint
selection, and training rules on three core datasets; Instacart
supplies a supplementary architecture comparison. The new
representative comparisons and independent evaluation are pending.
This draft therefore states hypotheses and protocol, without a
superiority claim.
\keywords{Event sequences \and Quantity forecasting \and
History correction \and Temporal point processes}
\end{abstract}

\section{Introduction}
Irregular demand sequences require both next-event quantity and
recorded time to be predicted. Recurrent and attention-based encoders
already condition these predictions on event history. We ask whether
an additional explicit local and recurrent path improves quantity
prediction within such an encoder. A full correction is compared
with its removal and with a current-position control of equal active
parameter count. The latter retains shared causal attention, so it
remains history-aware.

A secondary question asks whether separate time- and quantity-query
normalization, followed by fixed score combination, improves
prototype retrieval. We compare it with top-four averaging,
similarity weighting, and a general query under the same information
and extra-parameter budget. The prototype bank is learned on Train
and stays static during evaluation; it is not an online store of
observed events.

A completed MLP8 validation study motivates these questions but
does not prove them for the representative CNN+GRU54 architecture.
Its matched history-on/off pairs improved overall quantity RMSE in
18 of 18 cases, yet this comparison also changed capacity and
some quantity-selected checkpoints worsened timing likelihood.
The new study retains unfavorable quantity-tail, time, and cost
outcomes. It uses one representative model specification across
the core datasets rather than choosing a different winner for
each dataset.

\section{Related Work}
RMTPP summarizes event history in recurrent hidden vectors
\cite{du2016rmtpp}; deep renewal models predict demand size
and interdemand time from recurrent representations
\cite{turkmen2021renewal}. Attention-based A-NHP models
irregular event histories with temporal attention and residual
layers \cite{yang2022anhp}. S2P2 combines continuous-time
dynamics with event impulses and nonlinear residual
transformations \cite{chang2025s2p2}. History, residual
connections, and joint time--quantity modeling are thus
established ideas. An event quantity can itself be regarded
as a mark even without a categorical event-type ID.

Persistent-memory attention adds learned, input-independent
vectors \cite{sukhbaatar2019memory}, and product-key memory
selects learned slots sparsely \cite{lample2019product}.
The present study does not claim novelty for top-k retrieval
or learned static slots. Its retrieval hypothesis concerns
the specified split-query constraint under matched budgets.

For a later external comparison, the inspected S2P2 adapter
projects observed numerical gaps and quantities into the
S2P2 history encoder, reads its state immediately after the
last observed event, and attaches the shared prediction heads.
It is a common-head history-encoder adaptation, not a
reproduction of S2P2's native marked likelihood. Historical
external runs require separate source, population, Runtime,
head, and selection audits before they can substantiate a
claim about the representative architecture.

\section{Prediction Task and Shared Heads}
Let $\mathcal H_j=\{(\Delta_i,q_i)\}_{i\le j}$ be the observed
prefix. Position $j$ is the last observed event, and $j+1$
the withheld target. Its quantity and gap do not enter the
TitanTPP encoder readout. The input embedding is
\begin{equation}
 e_i=W_{\mathrm{in}}
 [\log(1+\Delta_i),\log(1+q_i)]^\top+b_{\mathrm{in}}+p_i .
 \label{eq:embedding}
\end{equation}
Here $p_i$ is learned positional information. These features
are not Train-z-scored. Two pre-normalized causal attention
blocks have width 64, four heads, feedforward width 128,
dropout 0.1, and 16 learned persistent attention tokens
per block. The persistent tokens differ from the 64 prototype
vectors applied after the encoder.

For final representation $z_j$, the quantity head produces
\begin{equation}
 \widehat\ell_{j+1}=\softplus(w_q^\top z_j+b_q),\qquad
 \widehat q_{j+1}=\exp(\widehat\ell_{j+1})-1 .
 \label{eq:quantity}
\end{equation}
Training compares $\widehat\ell_{j+1}$ with
$\log(1+q_{j+1})$ by squared error. Raw-unit evaluation
uses Eq.~\ref{eq:quantity}, with no later inverse Train
standardization.

The timing head predicts location $\mu_j$ and scale
$\sigma_j=10^{-3}+\softplus(w_\sigma^\top z_j+b_\sigma)$.
A latent duration $T=s\exp(\mu_j+\sigma_j\epsilon)$,
$\epsilon\sim\mathcal N(0,1)$, is recorded as the positive
integer $D=\max(1,\operatorname{round}(T))$.
For standard normal CDF $\Phi$, its probability mass is
\begin{align}
 p_j(1)&=\Phi\left(\frac{\log(1.5/s)-\mu_j}{\sigma_j}\right),
 \label{eq:p1}\\
 p_j(d)&=\Phi\left(\frac{\log((d+0.5)/s)-\mu_j}{\sigma_j}\right)
 -\Phi\left(\frac{\log((d-0.5)/s)-\mu_j}{\sigma_j}\right),
 \quad d\ge2. \label{eq:pd}
\end{align}
The core datasets have no upper top-code. Instacart
top-codes at 30 and uses the upper survival mass.
The target loss adds $-\log p_j(D_{j+1})$ to the log-quantity
squared error, with unit quantity coefficient and no
tail-weighted loss.

\section{Explicit History Correction}
The first block projects normalized inputs into queries,
keys, and values. For each $X\in\{Q,K,V\}$, it applies
\begin{equation}
 \widetilde X_i=X_i+
 \sum_{\ell=0}^{2}a_{i\ell}\kappa^X_\ell
 \odot X_{s_\ell(i)} .
 \label{eq:local}
\end{equation}
Here $s_0(i)=i$; $s_1(i)$ and $s_2(i)$ select the preceding
two observed positions in the current segment.
Eligibility $a_{i\ell}$ requires at least $\ell+1$
observed positions. Each depthwise kernel has 64
bias-free coefficients. Persistent key and value tokens
enter after this local operation. Let $b_i$ be the
complete first attention block output, including its
feedforward residual.

A width-54 recurrent path connects the blocks:
\begin{equation}
 u_i=A b_i,\quad r_i=\operatorname{GRU}(u_i,r_{i-1}),
 \quad c_i=\mathbf1[n_i\ge2]B r_i,\quad
 \widetilde b_i=b_i+c_i .
 \label{eq:gru}
\end{equation}
The matrices map $64\!\to\!54$ and $54\!\to\!64$.
The count $n_i$ includes the current observed position
in its segment; recurrent state begins at zero at
each segment boundary. The second block receives
$\widetilde b_i$. Thus the residual contains the
recurrent state updated by the current observation.
Padding contributes no correction; a valid withheld
position resets local eligibility and recurrent state.
This reset does not introduce an additional segment
mask into the shared causal attention.

The full path contains $3\cdot3\cdot64=576$ local
coefficients and 24,732 recurrent-path parameters:
25,308 active correction parameters in total.
\emph{Disabled} bypasses and freezes the path while
preserving both attention blocks, persistent tokens,
mean prototype retrieval, and prediction heads.
\emph{Current-only} retains the same correction
tensors. Its three local source values all come
from the current position:
\[X_i,\quad\tanh(X_i),\quad\operatorname{GELU}(X_i).\]
Its recurrent state starts at zero independently at
each position and processes the current $u_i$ twice,
with no state carried between positions.
This allows its recurrent weights to participate.
The first-block $b_i$ already contains prefix context,
whereas the local $X_i$ is computed before the first
attention operation. The equal parameter count does
not equalize functional capacity, depth, or cost.
A positive comparison would support the combined
three-position and recurrent correction path; it
would not isolate lag one, CNN, or GRU alone.

\section{Static Prototype Retrieval}
Let $M=\{m_k\}_{k=1}^{64}$ be train-learned
vectors of width 64. Define normalization by
\[\Nrm(v)=\frac{v}{\max(\|v\|_2,\varepsilon)}.\]
The normalized slot supplies its key; the same
unnormalized slot supplies its value. For second-block
output $h_j$, mean retrieval scores each slot by
\[s_{jk}=\Nrm(h_j)^\top\Nrm(m_k).\]
It selects the four highest-scoring slots $I_j$
and returns
\begin{equation}
 z_j^{\mathrm{mean}}=
 h_j+\tfrac14\sum_{k\in I_j}m_k .
 \label{eq:mean}
\end{equation}
Weighted retrieval replaces $\tfrac14$ with softmax
weights over those same four scores. Neither arm adds
query parameters.

The learned-query arms summarize the most recent
32 observed positions in the current segment, including
the current event. Each of time and quantity contributes
four features: last standardized value, mean,
population standard deviation, and last-minus-previous
difference. Positive time gaps use $\log(\Delta_i/s)$;
quantities use $\log(1+q_i)$, including zero quantities.
Train statistics standardize these points. Missing
features are zero. Write the views as
$\phi_j^t,\phi_j^q\in\mathbb R^4$.

The fixed arm learns
$W_t,W_q\in\mathbb R^{4\times64}$ and scores
\begin{align}
 q_j^t&=h_j+\phi_j^tW_t,\quad
 q_j^q=h_j+\phi_j^qW_q,\nonumber\\
 s_{jk}^{\mathrm{fixed}}
 &=\tfrac12\Nrm(q_j^t)^\top\Nrm(m_k)
 +\tfrac12\Nrm(q_j^q)^\top\Nrm(m_k).
 \label{eq:fixed}
\end{align}
The general arm learns $W\in\mathbb R^{8\times64}$:
\begin{equation}
 q_j=h_j+\tfrac12[\phi_j^t,\phi_j^q]W,\quad
 s_{jk}^{\mathrm{general}}=
 \Nrm(q_j)^\top\Nrm(m_k).
 \label{eq:general}
\end{equation}
Both add exactly 512 active parameters, select
four slots, and use a temperature-one softmax
for value aggregation. An earlier 521-parameter
general-query pilot is not this matched control.
The fixed score also equals one-half of the dot
product of the concatenated separately normalized
queries and duplicated normalized key. This
algebraic relation narrows the hypothesis to the
effect of the normalization and query constraint;
the labels "time" and "quantity" do not establish
semantic disentanglement. Every arm is fitted end
to end, so fitted encoder and slot values can differ.

\section{Experimental Design}
\subsection{Datasets, stages, and training}
The core study covers Taxi, Intermittent, and
RAF Spare Parts. Instacart is supplementary.
Table~\ref{tab:data} records frozen eligible
target populations. The legacy loader field
named \texttt{lookback\_weeks} counts each dataset's
native discrete time bins. Its inclusive window
ends at the last observed event; the next target
is appended before left-truncation to the
sequence cap. The cap includes the target.
Large-quantity thresholds were fixed from Train.

\begin{table}[t]
\caption{Frozen task populations and windows.
"Large" is strictly above the Train boundary.}
\label{tab:data}
\centering\small
\begin{tabular}{lrrrll}
\toprule
Dataset & Train & Validation & Cap & Lookback & Large\\
\midrule
Taxi & 38,393 & 8,268 & 256 & 168 hourly bins & $q>3449$\\
Intermittent & 393,824 & 86,285 & 256 &
520 weekly bins & $q>187$\\
RAF Spare Parts & 25,779 & 6,690 & 84 &
84 monthly bins & $q>200$\\
Instacart & 1,991,192 & 503,733 & 64 &
52 daily bins & $q>35$\\
\bottomrule
\end{tabular}
\end{table}

The campaign specifies 63 new fits in
three sequential stages (Table~\ref{tab:design}).
Instacart compares widths 8 and 12 with CNN+GRU54
across seeds 42, 52, and 62: nine fresh fits.
Historical width-16 checkpoints are reference
material, not new paired runs. The main correction
study trains full, disabled, and current-only
variants across three core datasets and three
seeds: 27 fresh fits. Retrieval trains weighted,
fixed, and general512 across the same datasets
and seeds: another 27. It reuses the nine new
full-correction fits as mean retrieval, exactly
once each. Same-dataset and same-seed arms
share a GPU and Runtime. Host timing is reported
separately.

\begin{table}[t]
\caption{Predeclared comparison. Status is the
writing-time record, not a live GPU observation.}
\label{tab:design}
\centering\small
\begin{tabular}{p{0.21\textwidth}p{0.46\textwidth}p{0.2\textwidth}}
\toprule
Stage & Arms and role & New fits/status\\
\midrule
Instacart & Width 8/12/CNN+GRU54;
architecture context & 9; in progress\\
Correction & Full/disabled/current-only;
primary claim & 27; pending\\
Retrieval & Weighted/fixed/general512;
full history reused as mean & 27; pending\\
\bottomrule
\end{tabular}
\end{table}

All fits use batch size 128, AdamW with learning
rate 0.001 and weight decay 0.01, gradient clipping
at 1, minimum 40 and maximum 300 epochs, and
patience 40. There is no warm start or automatic
retry. Checkpoint selection retains the first
strictly improving finite minimum of overall
raw-quantity RMSE on full Validation; ties keep
the earlier epoch. Neither a tail score nor
time likelihood changes that choice. A resource
timeout is not a completed scientific fit.

\subsection{Metrics, cost, and decisions}
At the quantity-selected checkpoint, report quantity
RMSE, quantity MAE, and recorded-time negative
log-likelihood both on all Validation targets and
on the large-quantity subset. All six values refer
to one selected epoch. Preserve each seed and
report three-seed means, sample standard
deviations, adverse cases, selected and final
saved epochs, fit time, and measured cost.
Recorded-time scores from different units
are not pooled.

The correction hypothesis receives developmental
support only if full has strictly lower overall
raw-quantity RMSE than \emph{both} controls for
all three seeds in the same two or more core
datasets. The retrieval hypothesis requires
fixed to beat mean, weighted, and general512
by the corresponding criterion. It also requires
cost eligibility on that dataset's GPU for
both lengths 64 and 256. Under batch 128,
five warm-up and twenty measured iterations,
relative to mean, inference time must rise
no more than 20\%, training-step time no
more than 30\%, and peak allocation no more
than 20\%. Technically successful but
over-budget arms remain in the report.

The supplementary Instacart study compares
architectures; it does not directly test either
correction or retrieval contribution. A result
below the strongest comparator would constrain
a universal superiority claim, while the
predeclared core comparisons would still decide
the two mechanism questions.

\section{Results and Interpretation}
\noindent\textbf{Writing-time status.}
The MLP8 Validation study and a limited parity
check between CNN+GRU54 and mean retrieval
are completed background evidence. The
representative correction and fixed-retrieval
comparisons are pending. Tables~\ref{tab:results}
and~\ref{tab:cost} contain \emph{pending}
entries so that no new performance result is
invented. The final version must fill them
from verified original checkpoints and report
every seed.

\begin{table}[t]
\caption{Reserved Validation comparison. All
metrics must come from each fit's quantity-selected
checkpoint, with per-seed results in the final version.}
\label{tab:results}
\centering\small
\begin{tabular}{@{}p{0.35\textwidth}p{0.22\textwidth}p{0.31\textwidth}@{}}
\toprule
Comparison & Overall quantity &
Large quantity and time\\
\midrule
Full vs.\ correction controls &
pending & pending\\
Fixed vs.\ retrieval controls &
pending & pending\\
Instacart architectures &
pending & pending\\
\bottomrule
\end{tabular}
\end{table}

\begin{table}[t]
\caption{Reserved same-GPU cost audit. Each
length requires inference, step, and peak allocation.}
\label{tab:cost}
\centering\small
\begin{tabular}{@{}p{0.17\textwidth}p{0.34\textwidth}p{0.34\textwidth}@{}}
\toprule
Arm & Length 64: inference, step, peak &
Length 256: inference, step, peak\\
\midrule
Mean & pending & pending\\
Weighted & pending & pending\\
Fixed & pending & pending\\
General512 & pending & pending\\
\bottomrule
\end{tabular}
\end{table}

\section{Discussion and Limitations}
An improvement over correction removal alone
would leave active capacity as a plausible
explanation. The current-only control narrows
this uncertainty, although equal parameters
do not equalize the operations or computation.
Even a successful comparison would support
the specified local-and-recurrent package,
not a lag-one-only or CNN-only effect.
Likewise, fixed retrieval has an equivalent
concatenated dot-product expression, so
performance must be attributed narrowly
to the tested constraint and protocol.

Development Validation is not untouched
evaluation. Earlier observations informed
the fixed-retrieval hypothesis, and historical
test partitions have prior research exposure.
A new independent assessment needs a verified
unexposed period or dataset, a frozen model
and selector, and an access-history audit
before outcomes are read. The 63-fit
campaign does not refit external S2P2.
Any final external table must distinguish
common-head adaptations from native model
likelihoods and disclose weaker comparisons.

Quantity selection may worsen timing
likelihood or rare large-quantity errors.
We therefore retain the same-epoch six
metrics and measured costs without setting
a favorable time threshold after the fact.
The sign-consistency rules do not by
themselves establish statistical significance
or universal external superiority.

\section{Conclusion}
This manuscript draft specifies a primary
test of explicit local and recurrent
correction within an already history-aware
event encoder. A separate matched-budget
study tests split-query static-prototype
retrieval. The eventual empirical
contribution depends on verified
three-seed comparisons, timing and cost
tradeoffs, and later independent
evaluation. Pending experimental outcomes
cannot yet support a final performance
claim.

\begin{thebibliography}{6}
\bibitem{du2016rmtpp}
Du, N., Dai, H., Trivedi, R., Upadhyay, U.,
Gomez-Rodriguez, M., Song, L.:
Recurrent marked temporal point processes:
Embedding event history to vector.
In: KDD, pp. 1555--1564 (2016)
\bibitem{turkmen2021renewal}
T\"urkmen, A.C., Januschowski, T., Wang, Y., Cemgil, A.T.:
Forecasting intermittent and sparse time series:
A unified probabilistic framework via deep renewal processes.
PLOS ONE 16(11), e0259764 (2021)
\bibitem{yang2022anhp}
Yang, C., Mei, H., Eisner, J.:
Transformer embeddings of irregularly spaced events
and their participants. In: ICLR (2022)
\bibitem{chang2025s2p2}
Chang, Y., Boyd, A., Xiao, C., Kass-Hout, T.,
Bhatia, P., Smyth, P., Warrington, A.:
Deep continuous-time state-space models for marked
event sequences. In: NeurIPS 38 (2025)
\bibitem{sukhbaatar2019memory}
Sukhbaatar, S., Grave, E., Lample, G., J\'egou, H.,
Joulin, A.: Augmenting self-attention with
persistent memory. arXiv:1907.01470 (2019)
\bibitem{lample2019product}
Lample, G., Sablayrolles, A., Ranzato, M.,
Denoyer, L., J\'egou, H.:
Large memory layers with product keys.
In: NeurIPS (2019)
\end{thebibliography}
\end{document}
